\section{本章小结}
\label{sec:tr-summary}

本章开头的问题是：局部模型只在附近可靠，怎样同时决定往哪里走、走多远。信赖域的回答是把两件事合并成一个带半径约束的子问题，半径表示当前愿意相信模型的范围。这样即使曲率矩阵不定，子问题也有解，负曲率不必被修正掉，反而可以把步推向边界。

子问题的全局解由一个平移过的线性方程刻画，但外层迭代并不需要它。Cauchy步给出充分模型下降的下限，狗腿法在$\bm{B}$正定时用两段折线逼近解的轨迹，截断共轭梯度只用矩阵向量乘积，并用边界和曲率两条规则处理信赖域带来的新情况。

下降比$\rho^{(t)}$把模型的预测和目标的实际下降作比较，决定接受这一步以及下一轮的半径。拒绝时点不动，模型沿用，只在更小的范围内重解。

收敛分析的条件和结论要分开记。目标下有界、梯度Lipschitz、模型有界、充分模型下降，这四条保证梯度趋于零，外层轮数是$O(\varepsilon^{-2})$，并不要求精确的Hessian。取精确Hessian并要求候选步利用负曲率，才有二阶必要条件的保证，而这仍不等于找到局部极小点，更不等于全局最优。

表~\ref{tab:tr-summary}把三种子问题解法放在一起。

\begin{BookInlineTable}
 \centering
 \begin{tabular}{llll}
  \toprule
  解法 & $\bm{B}$的要求 & 主要代价 & 负曲率\\
  \midrule
  Cauchy步 & 对称 & 一次乘积 & 只看负梯度方向\\
  狗腿法 & 正定 & 一次分解 & 不能处理\\
  截断共轭梯度 & 对称 & 若干次乘积 & 遇到时走到边界\\
  全局解 & 对称 & 数次分解 & 完整利用\\
  \bottomrule
 \end{tabular}
 \caption{子问题几种解法的对照。}
 \label{tab:tr-summary}
\end{BookInlineTable}

最后一列是第~\ref{chap:saddle-points}章的入口。模型处理负曲率的能力只是工具，梯度很小时所得驻点属于哪一种，怎样可靠地离开鞍点，要在那里进一步明确。
